\section{Sets of Numbers}

\begin{enumerate}
    \item [2.] $(-1,5] \cap [0, 8) = [0, 5]$
    
    \item [8.] $(-\infty, 5) \cup (5, \infty)$
    
    \item [20.] \leavevmode
    \begin{center}
        \includegraphics[width=0.5\textwidth]{figures/01_Relations_and_Functions/exercise_20.pdf}
    \end{center}
    
    \item [22.]
    \[ d = \sqrt{(-3-1)^2 + (5-2)^2} = \sqrt{16 + 9} = \sqrt{25} = 5\]
    \[M = (\frac{1-3}{2}, \frac{2+5}{2}) = (-1, 3.5)\]
    
    \item [30.]
    \[4 = \sqrt{(3-x)^2 + (2+1)^2} = \sqrt{9-6x+x^2 + 9} = \sqrt{x^2-6x+18}\]
    \[x^2-6x+18=16\]
    \[x^2-6x+2=0\]
    \[x_1 = 3 + \sqrt{7}, x_2 = 3-\sqrt{7}\]
    \[(3+\sqrt{7}, -1), (3-\sqrt{7}, -1)\]
    
    \item [34.] He's at $(-3,-4)$, at a distance of $5$ miles, after running he will be at $(4,-4)$
    
    \item [35.]
    \begin{enumerate}
        \item 
        \[d = \sqrt{(a-a)^2 + (y_1 - y_0)^2} = \sqrt{(y_1 - y_0)^2} = y_1 - y_0\]
        \item 
        \[d = \sqrt{(x_1 - x_0)^2 + (b-b)^2} = \sqrt{(x_1 - x_0)^2} = x_1 - x_0\]
        \item
        \[d = \sqrt{(x-x)^2 + (y-y)^2} = \sqrt{0} = 0\]
    \end{enumerate}
    
    \item [36.]
    \[M = (\frac{x_1 + x_2}{2}, \frac{y_1 + y_2}{2})\]
    \begin{align*}
    d_{PM} &= \sqrt{(\frac{x_1 + x_2}{2} - x_1)^2 + (\frac{y_1 + y_2}{2} - y_1)^2} \\
           &= \sqrt{(\frac{x_2-x_1}{2})^2 + (\frac{y_2-y_1}{2})^2} \\
           &= \sqrt{\frac{x_1^2-2x_1x_2+x_2^2}{4} + \frac{y_1^2-2y_1y_2+y_2^2}{4}} \\
    d_{QM} &= \sqrt{(\frac{x_1 + x_2}{2} - x_2)^2 + (\frac{y_1 + y_2}{2} - y_2)^2} \\
           &= \sqrt{(\frac{x_1-x_2}{2})^2 + (\frac{y_1-y_2}{2})^2} \\
           &= \sqrt{\frac{x_1^2-2x_1x_2+x_2^2}{4} + \frac{y_1^2-2y_1y_2+y_2^2}{4}}
    \end{align*}
    \[d_{PM} = d_{QM}\]

    \item [37.] To prove that the points are the vertices of a right triangle there must be an internal angle with a value of 90deg or prove the converse of the pythagorean theorem.
    \begin{enumerate}
        \item
        \[AB = \sqrt{13}, AC=2\sqrt{13}, BC=\sqrt{65}\]
        \[BC^2 = 65 = AB^2 + AC^2 = 13 + 52\]
        \item
        \[AB = 5\sqrt{2}, AC=5, BC=5\]
        \[AB^2 = 50 = AC^2 + BC^2 = 25 + 25\]
    \end{enumerate}
    
    \item [38.] To find the remaining point that makes the square, we must find a point that is at the same distance as the other points. We can create a vector from one of the sides of the square and just use it as a base to create the square.
    If we take $CB=(4,3)$, then applying that to A, we get $D=(1, 4)$ 

\end{enumerate}